\begin{textsample}{NEG Dim 1 – vem\_ed\_22 – Score -2.00 – vem\_ed\_22/t110.txt}  \label{ex:f1_neg_019}
Leitores Aos Osvaldo Succi Jr.
Coordenador PCIs Nesta edição, celebramos excelentes resultados conquistados em 2023, graças ao esforço e \textbf{dedicação} de todos os “PCIzeiros” – como carinhosamente chamamos os professores das Fatecs \textbf{engajados} nos PCIs – e ao apoio das equipes gestoras nas unidades.
Os principais resultados do ano estão \textbf{resumidos} na próxima página, entre eles: 109 PCIs envolvendo 2.450 estudantes de Fatecs, 2.505 de instituições internacionais, 91 professores de Fatecs, 78 estrangeiros.
Além dos números significativos, 2023 foi o ano em que os PCIs completaram 10 anos de \textbf{existência} e conquistaram um passo importante rumo à sua projeção internacional.
A equipe PCI / Cesu marcou presença em eventos internacionais, organizou o Simpósio Internacional IVES Cesu 2023 e a Jornada de PCIs e, pela primeira vez, ofereceu capacitações presenciais na DUOC UC (Chile).
Os PCIs desenvolvidos nas Fatecs ganham \textbf{maturidade} e expandem suas fronteiras, recebendo reconhecimento internacional.
É o caso dos trabalhos apresentados no Symbiosis Student´s Conclave, evento online organizado pelo Symbiosis Centre for Management Studies, universidade \textbf{indiana} que desenvolve PCIs com Fatecs desde 2021.
E do prêmio \textbf{concedido} pela Câmara de Comércio e Indústria dos BRICS ao PCI realizado entre Fatecs Itapetininga, Indaiatuba e São Caetano com a Durban University of Technology (África do Sul).
Tenho certeza de que, em 2024, conseguiremos trazer um pouco mais de cidadania global aos nossos alunos.
Boa leitura!

% matched lemmas: concedir, dedicação, engajados, existência, indiana, maturidade, resumido
\end{textsample}
